\documentclass[11pt,a4paper]{ctexart}

\usepackage[margin=2.2cm]{geometry}
\usepackage{amsmath}
\usepackage{array}
\usepackage{booktabs}
\usepackage{enumitem}
\usepackage{fancyhdr}
\usepackage{hyperref}
\usepackage{longtable}
\usepackage{pgfplots}
\usepackage{tabularx}
\usepackage{url}
\usepackage{xcolor}

\pgfplotsset{compat=1.18}

\hypersetup{
  colorlinks=true,
  linkcolor=blue!55!black,
  urlcolor=blue!65!black,
  pdftitle={周期产品数据分析：以 CF Industries 与氮肥周期为例},
  pdfauthor={stock-research}
}

\setlist{nosep,leftmargin=2em}
\setlength{\parindent}{2em}
\setlength{\parskip}{0.35em}
\setlength{\headheight}{14pt}
\renewcommand{\arraystretch}{1.25}
\newcolumntype{Y}{>{\raggedright\arraybackslash}X}
\newcolumntype{L}[1]{>{\raggedright\arraybackslash}p{#1}}

\pagestyle{fancy}
\fancyhf{}
\fancyhead[L]{周期产品数据分析：CF 与氮肥周期}
\fancyhead[R]{2026-07-16}
\fancyfoot[C]{\thepage}

\title{周期产品数据分析\\——以 CF Industries 与氮肥周期为例}
\author{stock-research}
\date{数据截面：2026-07-16\\文档版本：1.0}

\begin{document}

\maketitle

\begin{abstract}
本文以 CF Industries 和氮肥行业为例，介绍如何认识周期产品数据，并从外部商品价格、能源成本和季节性销量出发，逐步估算公司的实现售价、单位利润和季度毛利润。前半部分说明数据的来源、口径和限制，后半部分依次使用会计恒等式、单位经济分析、一元线性回归、等权分布滞后、季度季节均值和残差分解，展示周期产品研究中“先理解恒等式和线性传导，再研究残差”的基本方法。
\end{abstract}

\tableofcontents
\newpage

\section{数据集要回答什么问题}

CF Industries（下文简称 CF）是一家以氮肥为核心产品的生产商。研究 CF 的关键不是单独观察股价或某一个肥料价格，而是理解下面这条传导链：

\begin{verbatim}
农产品价格与农户利润
        ↓
种植面积、施肥意愿和采购节奏
        ↓
氨、尿素、UAN 等氮肥价格
        ↓
天然气成本与全球氮肥成本曲线
        ↓
CF 的实现售价、销量和单位利润
        ↓
毛利、EBITDA、现金流与估值
        ↓
CF 股价及市场预期
\end{verbatim}

因此，数据集围绕五类问题组织：

\begin{enumerate}
  \item 氮肥的市场售价处于什么水平？
  \item 天然气成本对生产利润形成多大压力？
  \item 玉米价格、种植面积和农户利润是否支持氮肥需求？
  \item 外部市场价格是否传导到了 CF 的实际售价和财务结果？
  \item CF 股价已经反映了多少周期变化？
\end{enumerate}

\section{数据内容总览}

\begin{center}
\small
\begin{tabularx}{\textwidth}{L{2.7cm}L{3.1cm}L{2.1cm}Y}
\toprule
数据类别 & 主要内容 & 典型频率 & 主要用途 \\
\midrule
股票市场 & CF 股价、成交量、收益和动量 & 交易日/月 & 观察市场预期和股票确认信号 \\
天然气 & Henry Hub 现货价格 & 日 & 观察北美氮肥的主要原料成本 \\
氮肥价格 & 全球尿素、美国氨、尿素和 UAN & 月/周 & 观察产品售价和区域市场强弱 \\
农业需求 & 玉米、大豆价格、种植面积、成本与回报 & 日/年内发布/年 & 观察农户支付能力和潜在施肥需求 \\
CF 财务 & 收入、毛利、库存、现金流、债务等 & 季度 & 观察周期向会计结果的传导 \\
CF 经营 & 分产品售价、销量、毛利/吨、实际气价 & 季度 & 解释 CF 的产品经济性 \\
派生指标 & 尿素气价差、CF 实现气价差、价格篮子等 & 日/月/季 & 用统一口径概括氮肥周期 \\
\bottomrule
\end{tabularx}
\normalsize
\end{center}

\section{市场和行业数据}

\subsection{CF 股票行情}

股票行情包含 CF 的开盘价、最高价、最低价、收盘价、复权收盘价和成交量，当前覆盖 2013-01-02 至 2026-07-15，共 3,403 个交易日。

\textbf{主要用途：}

\begin{itemize}
  \item 计算股票日收益、月收益和未来区间收益；
  \item 计算 20、60、120 个交易日的动量；
  \item 计算短期和中期波动率；
  \item 与氮肥价差、财务数据和估值数据比较。
\end{itemize}

\textbf{理解要点：}复权价格适合计算持有期收益；未复权收盘价更适合与当时市值、企业价值等截面数据配合。股票价格是市场对未来的综合预期，通常不能被当作基本面本身。

\subsection{Henry Hub 天然气}

Henry Hub 是美国天然气现货基准，单位为美元/MMBtu。数据来源于 EIA，当前覆盖 1997-01-07 至 2026-07-13，共 7,410 条日度记录。

天然气是氨生产的原料和能源，也是 CF 成本优势的核心来源之一。Henry Hub 下降通常有利于北美氮肥生产利润，但实际影响还取决于：

\begin{itemize}
  \item CF 的实际采购价格、套保和运输成本；
  \item 欧洲及其他边际产区的天然气价格；
  \item 氮肥售价是否同步下降；
  \item 库存成本和产品销售时点。
\end{itemize}

因此，Henry Hub 是重要成本代理，但不等同于 CF 披露的实际天然气成本。

\subsection{全球尿素价格}

全球尿素数据来自 World Bank Pink Sheet，单位为美元/公吨，频率为月度。当前历史覆盖 1960-01-01 至 2025-12-01，共 792 条记录。

这组数据的优势是历史长、口径相对稳定，适合观察全球氮肥的长周期位置。它的局限包括：

\begin{itemize}
  \item 月度频率不能反映周度或日度市场变化；
  \item 全球基准与北美内陆、NOLA 或 CF 实现价格之间存在地区基差；
  \item 截至本文数据日期，最新观测距离当前已有 227 天，近期值已经陈旧；
  \item 单一尿素价格不能完整代表氨、UAN 和其他产品。
\end{itemize}

因此，这组数据适合作为\textbf{长期全球锚}，不适合单独承担近期市场判断。

\subsection{USDA AMS 美国肥料价格}

AMS 数据来自 USDA Agricultural Marketing Service 的 Illinois Production Cost Report，当前覆盖 2022-09-22 至 2026-07-10，共 321 条产品报价。主要产品包括：

\begin{itemize}
  \item 无水氨（anhydrous ammonia）；
  \item 46\% 尿素（urea 46）；
  \item 28\% UAN；
  \item 32\% UAN。
\end{itemize}

价格单位为美元/短吨，通常记录最低价、最高价和平均价。数据频率大致为周度或双周。

AMS 数据比全球月度尿素更接近美国终端市场，可以帮助观察北美现货是否确认全球价格方向。但它并不是 CF 的公司实现价：AMS 是 Illinois 经销商报价，而 CF 的销售跨越不同地区、客户、运输方式和合同期限。

AMS 历史从 2022 年开始，覆盖不足一个完整长周期，因此更适合观察\textbf{当前区域市场和短期基差}，而不是独立训练长期模型。

\subsection{玉米和大豆期货}

农业价格包括玉米和大豆连续期货代理，单位为美分/蒲式耳。当前覆盖 2013-01-02 至 2026-07-16，两类作物合计 6,820 条记录。

玉米是北美氮肥需求中最重要的作物之一，单位面积氮肥投入通常高于大豆。因此：

\begin{itemize}
  \item 玉米价格上升通常改善玉米种植收益和施肥支付能力；
  \item 玉米相对大豆更有吸引力时，可能推动更多面积转向玉米；
  \item 大豆价格仍然重要，因为玉米和大豆存在种植面积竞争；
  \item 农产品价格下跌不一定立即导致氮肥需求下降，库存、保险、远期销售和成本结构都会造成滞后。
\end{itemize}

连续期货存在合约换月和期限结构影响，价格跳变不一定完全来自现货基本面。

\subsection{玉米和大豆种植面积}

种植面积数据来自 USDA NASS，单位为英亩，当前覆盖 2013 年至 2026 年，共 106 条不同作物和发布口径的记录。

面积数据用于估计潜在氮肥需求规模。使用时必须注意：

\begin{itemize}
  \item 意向面积、实际播种面积和最终修订面积不是同一个概念；
  \item 一年中的多次发布是对同一作物年度的更新，不能当作完全独立样本；
  \item 面积只决定需求基数，单位面积施肥量、天气和播种进度同样重要；
  \item 部分历史记录的可用日期不一定精确还原原始发布日期，严格历史回测时需谨慎。
\end{itemize}

\subsection{农业成本与回报}

玉米和大豆成本与回报数据来自 USDA ERS。玉米覆盖 1996---2025 年，共 5,665 条分地区、分成本项目记录；大豆覆盖 1997---2025 年，共 6,026 条记录。

行数较多是因为每个年度包含多个地区、成本类别和收入项目，并不代表有数千个独立年度样本。常见项目包括种子、肥料、化学品、燃料、土地、人工、总成本和总收入。

这类数据适合构造农户利润情景、解释长期支付能力，但频率低、修订多，不适合像日度价格一样直接进入高频分析。

\section{CF 公司数据}

\subsection{财务数据}

CF 财务数据来自 SEC XBRL Company Facts，当前原始记录覆盖 2006-12-31 至 2026-05-04，共 27,212 条不同概念、期间和申报记录。

研究中主要使用：

\begin{itemize}
  \item 营业收入和毛利润；
  \item 净利润和营业利润；
  \item 折旧与摊销；
  \item 经营现金流和资本开支；
  \item 库存、现金和长期债务；
  \item 流通股数。
\end{itemize}

财务数据存在几个常见口径问题：

\begin{itemize}
  \item 同一经济含义可能使用不同 XBRL concept；
  \item 10-Q 中的现金流和部分利润指标可能按年初至今累计披露；
  \item 后续申报可能重述历史期间；
  \item 财务期末日早于实际发布日期，研究时不能在期末日提前使用数据。
\end{itemize}

\subsection{分产品经营数据}

分产品经营数据来自 CF 的季度和年度经营披露，当前覆盖 2015-06-30 至 2026-03-31，共 834 条产品和指标记录。

主要产品包括：

\begin{itemize}
  \item 氨；
  \item 颗粒尿素；
  \item UAN；
  \item 硝酸铵；
  \item 其他产品。
\end{itemize}

主要指标包括平均售价、销量、毛利/吨、全产品产量和实际天然气成本。常见单位为美元/短吨、千短吨和美元/MMBtu。

这组数据是连接外部氮肥市场与 CF 财务结果的关键：外部报价反映市场环境，CF 实现售价反映合同、地区、运输和销售时点后的公司结果。

\section{研究用汇总数据表}

为了便于研究，不同频率的数据被整理成若干张具有固定观察粒度的数据表。下面只介绍它们包含什么以及应如何理解。

\subsection{日度研究表}

日度研究表以 CF 的每个交易日为一行，当前有 3,403 行、62 个字段，覆盖 2013-01-02 至 2026-07-15。

字段主要分为：

\begin{itemize}
  \item CF 行情、收益、动量和波动率；
  \item 最新可见的 Henry Hub、尿素、玉米、大豆和 AMS 价格；
  \item 春季和秋季施肥季节标识；
  \item 最新已经披露的 CF 财务和经营指标；
  \item 每项低频数据对应的原始观察日期和公开日期。
\end{itemize}

需要特别注意：季度财务、年度面积或月度价格在日度表中会重复出现在多个交易日。这种展开便于按日期查询，但重复的行不构成新的独立经济信息。

\subsection{季度经营表}

季度经营表以 CF 的每个财务季度为一行，当前有 42 行、84 个字段，覆盖 2015-06-30 至 2026-03-31。

它把分产品经营、公司财务、季度市场均价和季度末估值放在同一行中，适合分析：

\begin{itemize}
  \item 氮肥价格变化如何传导至 CF 实现售价；
  \item 天然气价格如何传导至 CF 实际气价和毛利；
  \item 销量、检修和产品结构如何影响季度利润；
  \item 毛利、EBITDA、现金流和估值如何随周期变化。
\end{itemize}

季度表只有约 42 个观察值，适合关系校准和经营桥分析，不适合同时使用大量变量训练复杂模型。

\subsection{核心月度表}

核心月度表把最重要的长期变量压缩到月度粒度，当前有 163 个月，其中 115 个月具备全部核心数据。

\begin{center}
\small
\begin{tabularx}{\textwidth}{L{4.4cm}Y}
\toprule
指标 & 含义 \\
\midrule
全球尿素气价差月均值 & 全球尿素价格扣除估算天然气成本后的外部利润环境 \\
Henry Hub 月均值 & 北美氮肥的主要原料成本 \\
玉米价格月均值 & 农业需求和农户支付能力代理 \\
CF 六个月股价动量 & 股票市场对周期方向的确认 \\
最新 CF 实现组合气价差 & 最新披露的公司产品售价与实际气价关系 \\
最新 CF 毛利率变化 & 公司盈利方向是否改善 \\
\bottomrule
\end{tabularx}
\normalsize
\end{center}

这张表适合做周期方向、领先滞后和简单统计关系分析。115 个完整月份仍然是小样本，应限制变量和滞后项数量。

\subsection{核心季度表}

核心季度表保留五项最能概括公司经营周期的指标，当前有 42 个季度，其中 40 个季度数据完整。

\begin{center}
\small
\begin{tabularx}{\textwidth}{L{4.4cm}Y}
\toprule
指标 & 含义 \\
\midrule
全球尿素气价差季度均值 & 外部氮肥周期的季度环境 \\
CF 实现组合气价差 & 公司实际售价扣除估算气价成本后的实现环境 \\
CF 总销量 & 需求、检修、运行率和库存共同作用后的销售结果 \\
CF 实际天然气成本 & 公司披露的实际气价，而非 Henry Hub 现货价 \\
CF 毛利率变化 & 盈利方向的一阶变化 \\
\bottomrule
\end{tabularx}
\normalsize
\end{center}

外部尿素气价差与 CF 实现组合气价差的样本内相关性约为 0.966。两者经济含义不同：前者描述外部市场，后者描述公司实际传导；但在只有约 40 个季度的统计分析中，不宜不加选择地同时使用。

\subsection{战术观察表}

战术观察表保存更适合当前场景判断的数据，包括：

\begin{itemize}
  \item AMS 氨、尿素和 UAN 价格；
  \item 美国尿素理论气价差；
  \item 固定权重的 CF 氮肥价格篮子；
  \item 玉米和大豆种植面积；
  \item 春季和秋季施肥窗口；
  \item 分产品其他成本和地区基差残差。
\end{itemize}

这张表适合回答“北美现货是否确认全球方向”“种植面积是否构成需求冲击”“公司成本残差是否出现异常”等问题。由于 AMS 历史较短、面积频率较低，它不适合作为长期基础样本。

\section{氮肥经济性派生指标}

\subsection{单位标准化}

数据中同时出现公吨、短吨、产品吨和含氮吨。主要换算关系为：

\begin{align*}
1\ \text{公吨} &= 1.10231131\ \text{短吨}, \\
P_{\text{美元/短吨}} &= 0.90718474\times P_{\text{美元/公吨}}.
\end{align*}

实际数据处理中，1 美元/公吨对应约 0.90718474 美元/短吨。产品含氮比例不同：氨约 82\%，尿素约 46\%，UAN 32 约 32\%，硝酸铵约 34\%。因此，同样的“每吨产品价格”不代表同样的“每吨氮价格”。

\subsection{理论现金价差}

理论现金价差用于粗略描述产品售价扣除天然气消耗后的空间：

\begin{equation*}
\text{理论现金价差}
=\text{标准化产品价格}
-\text{天然气价格}\times\text{单位产品气耗}.
\end{equation*}

基准气耗假设如下：

\begin{center}
\begin{tabular}{l r}
\toprule
产品 & 基准气耗（MMBtu/短吨） \\
\midrule
氨 & 28.0 \\
颗粒尿素 & 19.0 \\
UAN & 13.0 \\
硝酸铵 & 14.0 \\
\bottomrule
\end{tabular}
\end{center}

低、基准和高气耗情景分别使用基准值的 0.85、1.00 和 1.15 倍。

这些气耗并非每座装置的工程参数，而是以 CF 公开披露的组合数据校准。CF 2023 年约 3.41 亿 MMBtu 气耗和 1,910 万产品吨隐含约 17.85 MMBtu/产品吨；当前产品组合假设回算约 17.95，误差低于 1\%。

\subsection{CF 实现气价差}

CF 实现气价差使用 CF 披露的产品平均售价和实际天然气成本：

\begin{equation*}
\text{CF 实现气价差}
=\text{CF 产品实现售价}
-\text{CF 实际气价}\times\text{气耗假设}.
\end{equation*}

它比外部市场价格更接近公司结果，但仍不等于会计毛利。以下项目尚未被完整扣除：

\begin{itemize}
  \item 运费、仓储和地区基差；
  \item 套保影响和库存计价时点；
  \item 外购产品；
  \item 检修、人工和其他制造成本；
  \item 不同产品和地区的销售组合。
\end{itemize}

实现气价差与披露毛利/吨之间的差额被保留为“其他成本与基差残差”，用于解释上述未建模因素。

\subsection{CF 氮肥价格篮子}

价格篮子把美国氨、尿素和 UAN 价格标准化后按固定权重合成：

\begin{center}
\begin{tabular}{l r}
\toprule
组成 & 权重 \\
\midrule
AMS 氨 & 23.95\% \\
AMS 尿素 & 31.47\% \\
AMS UAN 32 & 44.58\% \\
\bottomrule
\end{tabular}
\end{center}

权重参考 CF 2019---2021 年产品销量结构，价格基期为 2024 年。该篮子用于观察美国氮肥组合价格方向，不代表 CF 的实际收入结构会永久保持不变。

\section{从行业数据到公司利润：经营传导的统计估算}

\subsection{分析框架：从会计恒等式到残差}

周期产品的经营结果往往由少数可解释的数量关系决定。售价来自外部商品价格和公司基差，主要原料成本接近“价格乘以耗用量”，单位利润是售价减成本，总利润则是单位利润乘销量。只要公司没有跨过停产、产能瓶颈或贸易制度突变等边界，这些关系在一个周期内通常接近线性。

因此，在使用多变量预测方法之前，可以先把问题分成三个层次：

\begin{enumerate}
  \item \textbf{恒等式：}例如分产品销量乘毛利/吨应当能够还原公司毛利润；
  \item \textbf{稳定参数：}例如单位产品气耗、季度销量基数和其他单位成本；
  \item \textbf{简单敏感度：}例如外部尿素每变化 1 美元/吨，CF 尿素实现价平均变化多少。
\end{enumerate}

只有这些简单关系留下的残差，才需要继续寻找合同滞后、地区基差、库存、套保、检修或结构变化等解释。

以下估算使用当前约 40---42 个 CF 季度观察值。所有系数均为样本内描述性结果，用于建立量级感和经营桥，不应直接解释为因果关系或未来预测保证。

\subsection{第一步：由外部价格估算 CF 实现售价}

先把 World Bank 尿素从美元/公吨换算为美元/短吨，记为 $X_t$。最简单的售价关系为：

\begin{equation*}
\widehat{P}_{p,t}=\alpha_p+\beta_p X_t,
\end{equation*}

其中 $\widehat{P}_{p,t}$ 是 CF 产品 $p$ 的季度实现售价。逐产品的一元估算结果如下：

\begin{center}
\small
\begin{tabularx}{\textwidth}{L{2.5cm}Y r r r r}
\toprule
产品 & 一元线性估算式 & 样本 & $R^2$ & MAE & 变化相关 \\
\midrule
颗粒尿素 & $87.7+0.849X_t$ & 41 & 0.931 & 27.9 & 0.770 \\
UAN & $40.0+0.681X_t$ & 41 & 0.804 & 37.2 & 0.468 \\
氨 & $89.1+1.108X_t$ & 41 & 0.735 & 73.5 & 0.383 \\
硝酸铵 & $99.7+0.589X_t$ & 40 & 0.826 & 26.0 & 0.599 \\
\bottomrule
\end{tabularx}
\normalsize
\end{center}

MAE 为平均绝对误差，单位是美元/短吨。“变化相关”比较的是相邻季度变化，而不是价格水平。

颗粒尿素的关系最直接：

\begin{equation*}
\widehat{P}_{\text{CF尿素},t}
=87.7+0.849\times P_{\text{全球尿素},t}.
\end{equation*}

它能够解释约 93\% 的历史价格水平变化，平均误差约 28 美元/短吨。UAN、氨和硝酸铵也跟随全球尿素周期，但截距、斜率和误差不同，反映了产品结构、地区和合同口径差异。

水平拟合普遍好于季度变化拟合。这意味着线性关系适合估算“周期大致处于什么利润水平”，但单个季度的转折时点仍会受到合同滞后、库存和基差影响。特别是氨的价格变化相关性只有约 0.38，不能仅凭全球尿素预测其季度变化。

\subsection{第二步：由 Henry Hub 估算 CF 实际气价}

CF 披露的实际天然气成本与 Henry Hub 当季均价存在明显线性关系：

\begin{equation*}
\widehat{G}^{\text{CF}}_t
=0.29+1.00\times HH_t.
\end{equation*}

这一关系的 $R^2$ 约为 0.755，MAE 约为 0.47 美元/MMBtu。如果把当季和上季 Henry Hub 取平均：

\begin{equation*}
\overline{HH}_t=\frac{HH_t+HH_{t-1}}{2},
\end{equation*}

引入一期等权分布滞后后，关系改善为：

\begin{equation*}
\widehat{G}^{\text{CF}}_t
=-0.21+1.17\times\overline{HH}_t.
\end{equation*}

$R^2$ 提高到约 0.898，MAE 降至约 0.30 美元/MMBtu。这个结果符合经营直觉：公司实际气价不是某一天的现货价，而是采购、运输、库存和套保共同形成的平滑价格。

这也说明，很多看似“非线性”的偏差实际上可能只是时间错位。先尝试当期、滞后一期和简单移动平均，通常比直接增加复杂函数更容易解释。

\subsection{第三步：用售价减成本估算单位利润}

产品 $p$ 的单位毛利可以先写成：

\begin{equation*}
\widehat{M}_{p,t}
=\widehat{P}_{p,t}
-g_p\widehat{G}^{\text{CF}}_t
-C_p,
\end{equation*}

其中 $g_p$ 是单位产品气耗，$C_p$ 是未单独披露的其他单位成本和基差。利用历史实现气价差与披露毛利/吨的差额，可以得到下面的固定成本量级：

\begin{center}
\small
\begin{tabularx}{\textwidth}{L{2.8cm}r r r Y}
\toprule
产品 & $C_p$ 均值 & 固定值 MAE & 自由拟合 $R^2$ & 解释 \\
\midrule
氨 & 201.7 & 38.3 & 0.911 & 成本水平较高，季度波动也较大 \\
颗粒尿素 & 140.9 & 20.0 & 0.967 & 线性单位经济关系较稳定 \\
UAN & 116.2 & 12.6 & 0.982 & 四类产品中最稳定 \\
硝酸铵 & 189.8 & 45.6 & 0.380 & 固定成本近似较弱，需要单独解释 \\
\bottomrule
\end{tabularx}
\normalsize
\end{center}

表中 $C_p$ 和 MAE 的单位均为美元/短吨。

例如，颗粒尿素可以先估算为：

\begin{equation*}
\widehat{M}_{\text{尿素},t}
=\widehat{P}_{\text{尿素},t}
-19\times\widehat{G}^{\text{CF}}_t
-141.
\end{equation*}

UAN 可以近似为：

\begin{equation*}
\widehat{M}_{\text{UAN},t}
=\widehat{P}_{\text{UAN},t}
-13\times\widehat{G}^{\text{CF}}_t
-116.
\end{equation*}

这里的 $C_p$ 不是严格意义上的固定制造成本，而是把运费、库存、检修、人工、外购产品、地区基差和产品结构暂时合并后的统计残差。它适合第一轮估算，但不应被理解为公司披露的成本科目。

颗粒尿素和 UAN 的线性解释力很高，表明“售价减天然气减稳定残差”已经能够解释大部分单位毛利变化。硝酸铵的 $R^2$ 只有约 0.38，是当前最值得继续拆解的产品；这可能来自地区、装置、会计口径或其他成本，而不一定意味着需要非线性算法。

\subsection{第四步：用季节均值估算销量}

CF 总销量更像一个受产能约束的稳定基数，而不是对农产品价格高度弹性的需求函数。将五类产品销量相加，并排除一个产品披露不完整的早期季度后，历史季度均值为：

\begin{center}
\begin{tabular}{l r r}
\toprule
季度 & 平均总销量（千短吨） & 完整季度数 \\
\midrule
第一季度 & 4,528 & 11 \\
第二季度 & 5,109 & 10 \\
第三季度 & 4,504 & 10 \\
第四季度 & 4,878 & 10 \\
\bottomrule
\end{tabular}
\end{center}

完整样本的总体均值约为 4,749 千短吨，标准差约为 422 千短吨，变异系数约 8.9\%。因此，第一版销量估算可以简单写成：

\begin{equation*}
\widehat{V}_t
=\text{季度季节均值}
+\Delta V_{\text{检修}}
+\Delta V_{\text{停产}}
+\Delta V_{\text{库存与销售时点}}.
\end{equation*}

农业需求变化首先影响的是行业价格和区域基差，而 CF 的短期销量还受到装置产能、计划检修和库存安排约束。因此，用玉米价格直接拟合 CF 总销量未必比“季节均值加显式事件调整”更可靠。

销量关系中真正的非线性主要来自离散事件：装置不是按连续比例缓慢减产，而可能因为检修、事故、天气或极端气价突然停产。对这类事件，加入单独的事件冲击项通常比拟合平滑曲线更直观。

\subsection{第五步：由单位利润和销量还原毛利润}

一旦得到分产品单位毛利和销量，毛利润就是一个近似恒等式：

\begin{equation*}
\widehat{GP}_t
=\sum_p \widehat{M}_{p,t}\widehat{V}_{p,t}.
\end{equation*}

当销量单位为千短吨、单位毛利为美元/短吨时，相乘后得到千美元；再除以 1,000 即为百万美元。

直接用 CF 已披露的分产品销量和毛利/吨重算公司毛利润，42 个季度的相关性约为 0.9998，中位绝对勾稽误差约为 0.59 百万美元。2018 年以后的平均绝对误差约为 0.57 百万美元。早期两个季度存在约 4,000 万美元的较大口径差异，需要在使用早期数据时单独核对。

这个勾稽结果说明：在“实现售价、实际气价、其他单位成本和销量”已经给定后，毛利润不需要另行拟合，直接按经营恒等式计算即可。真正需要估算的是恒等式右侧的少数经营输入。

\subsection{完整经营桥与残差分解}

把以上关系合并，可以形成一张低、中、高情景都能手算的 CF 经营桥：

\begin{enumerate}
  \item 输入全球尿素价格，按各产品的截距和斜率估算 CF 实现售价；
  \item 输入当季和上季 Henry Hub，估算 CF 实际天然气成本；
  \item 按产品气耗和历史其他单位成本估算毛利/吨；
  \item 以季度销量均值为基准，对检修、停产和库存进行显式调整；
  \item 逐产品计算“销量乘毛利/吨”，加总得到季度毛利润；
  \item 将实际结果与估算结果比较，把误差分解为售价基差、气价差、其他成本差和销量差。
\end{enumerate}

用符号表示，这张经营桥是：

\begin{align*}
\widehat{P}_{p,t} &= \alpha_p+\beta_pX_t, \\
\widehat{G}^{\text{CF}}_t &= a+b\frac{HH_t+HH_{t-1}}{2}, \\
\widehat{M}_{p,t} &= \widehat{P}_{p,t}-g_p\widehat{G}^{\text{CF}}_t-C_p, \\
\widehat{V}_{p,t} &= \overline{V}_{p,q}+\Delta V_{p,t}, \\
\widehat{GP}_t &= \sum_p\widehat{M}_{p,t}\widehat{V}_{p,t}.
\end{align*}

这不是为了追求最高样本内拟合度，而是为了让每一个误差都有明确含义：

\begin{itemize}
  \item 售价误差对应合同滞后、地区和产品基差；
  \item 气价误差对应采购、运输、库存和套保；
  \item 单位成本误差对应检修、人工、外购产品和会计口径；
  \item 销量误差对应装置运行、库存和销售时点。
\end{itemize}

如果某类残差只是围绕稳定均值波动，均值、季节项和线性敏感度已经能够概括主要关系；如果残差持续偏向一侧或在特定环境突然放大，才说明存在结构变化、遗漏变量或分段关系，值得进入更深入的统计分析。

\subsection{线性关系的边界}

现有数据支持“基础传导大体线性”的判断，但线性并不是无条件成立。以下情况可能造成分段或跳变：

\begin{itemize}
  \item 天然气或氮肥价格跨过海外装置停产与复产阈值；
  \item CF 装置检修、事故、极端天气或天然气供应中断；
  \item 制裁、出口限制、关税或航运中断改变地区基差；
  \item 经销商和农户在旺季前后集中补库或去库；
  \item CF 的产品、地区或合同结构发生明显变化；
  \item 会计披露范围或分产品口径改变。
\end{itemize}

这些因素多数可以先通过情景调整、事件标记或分阶段参数处理。只有当简单分段仍无法解释稳定残差时，复杂非线性方法才可能带来真正的增量价值。

\section{经营桥的滚动检验}

\subsection{实验问题与可证伪假设}

本实验并不检验经营桥能否解释已经公布的财报，而是检验：\textbf{在 CF 财报公布之前，经营桥能否比“直接沿用上一期实际值”更准确地估算当季经营结果。}

对每一个预测目标，定义朴素对照为预测时点最新已经披露的上一期实际值。可证伪假设为：

\begin{description}[style=nextline]
  \item[零假设] 经营桥的样本外误差不低于朴素对照，经营桥没有提供增量信息；
  \item[备择假设] 经营桥在相同季度上的样本外误差更低，并且方向判断更稳定。
\end{description}

实验同时评价中间经营变量和最终毛利润。如果最终毛利润改善，但售价、气价、销量和单位毛利都没有改善，就需要警惕误差偶然抵消；反过来，中间变量改善但毛利润没有改善，则说明误差在乘法和加总过程中被放大。

\subsection{预测目标与单位}

每个样本外季度最多产生 15 条预测：

\begin{center}
\small
\begin{tabularx}{\textwidth}{L{3.0cm}L{2.0cm}L{3.3cm}Y}
\toprule
任务组 & 目标数量 & 单位 & 目标内容 \\
\midrule
实际气价 & 1 & 美元/MMBtu & CF 全产品实际天然气成本 \\
实现售价 & 4 & 美元/短吨 & 氨、颗粒尿素、UAN、硝酸铵 \\
单位毛利 & 4 & 美元/短吨 & 上述四类产品的披露毛利/吨 \\
销量 & 5 & 千短吨 & 四类主要产品及其他产品 \\
毛利润 & 1 & 百万美元 & CF 公司季度毛利润 \\
\bottomrule
\end{tabularx}
\normalsize
\end{center}

实际气价和销量有 30 个有效样本外季度。由于 2026 年第一季度缺少当季全球尿素数据，依赖尿素价格的实现售价、单位毛利和毛利润有 29 个有效季度。评价时不以插值补造该季度，而是对经营桥和朴素对照同时使用共同有效样本。

\subsection{预测时点与信息边界}

一个季度的实验时间线为：

\begin{verbatim}
季度开始                    季度结束        市场数据就绪       形成预测       CF 财报公开
   |                            |                |                |               |
   +----------------------------+----------------+----------------+---------------+
                              T 日          约 T+14 日        T+15 日       T+31 至 T+56 日
\end{verbatim}

预测时点选择季度末后第 15 天，原因是：

\begin{itemize}
  \item Henry Hub 季度数据最晚通常在季度末后 1---3 天可用；
  \item World Bank 当季最后一个月的尿素数据通常在季度末后约 14 天可用；
  \item 当前样本中 CF 季度结果最早在季度末后 31 天公开；
  \item 因此，第 15 天可以使用完整的季度外部市场信息，但仍早于公司结果。
\end{itemize}

这里的实验是\textbf{季度末经营结果估算}，不是季度初预测。如果以后需要在季度开始前形成预测，必须另设信息截止点，不能与本实验结果混合比较。

\subsection{为什么不能用全样本系数评价效果}

前一节用全部历史季度计算截距、斜率和平均成本，目的是说明经济关系和数量级。这些全样本结果包含了后期数据，只能用于描述，不能证明当时真的可以估算出未来季度。

要检验经营桥是否具有实际增量，需要模拟每个历史时点的信息集合：在季度结束后选择一个固定日期，只使用当时已经公开的市场数据和公司历史结果，重新估计参数，再预测尚未披露的当季 CF 经营结果。

当前滚动检验采用以下规则：

\begin{itemize}
  \item 预测时点固定为季度末后第 15 天；
  \item CF 季度结果最早在季度末后 31 天公开，因此预测时尚不知道目标值；
  \item 每次只使用预测时点之前已经披露的公司季度；
  \item 初始训练样本至少为 12 个季度，此后采用扩展窗口；
  \item 外部尿素和 Henry Hub 使用截至预测时点已经发布的季度数据；
  \item 朴素对照为预测时点最新已披露的上一期实际值；
  \item 每条预测保存训练期起止、最后训练数据公开时间和当期公式参数；
  \item 评价指标包括 MAE、RMSE、偏差、WAPE、方向准确率和相对改善率。
\end{itemize}

这种设计检验的是“在财报发布前约两周，经营桥能否比简单沿用上一季度结果提供更多信息”，不是在季度开始前预测整个季度，也不是解释已经发布的财报。

\subsection{每个滚动季度的计算步骤}

设当前待预测季度为 $t$，预测截止时间为 $c_t$。每轮实验执行：

\begin{enumerate}
  \item 选择所有 $period<t$ 且公司结果公开时间不晚于 $c_t$ 的历史季度作为训练集；
  \item 如果训练季度少于 12 个，则不产生样本外预测；
  \item 在历史训练集上重新估计各产品售价的一元线性回归；
  \item 用当季与上季 Henry Hub 的等权平均，在历史训练集上重新估计 CF 实际气价；
  \item 用历史其他成本残差均值、固定气耗和预测售价/气价计算单位毛利；
  \item 用相同日历季度的历史销量均值估算分产品销量；
  \item 将分产品预测销量乘预测单位毛利后加总，得到公司毛利润；
  \item 同时记录上一期已披露实际值作为朴素对照；
  \item 财报公布后写入实际值，但不反向修改已经形成的预测。
\end{enumerate}

扩展窗口意味着训练集随时间增加，但不会删除早期季度。后续如研究结构变化，可以增加固定长度滚动窗口作为候选方法；它仍必须使用同一个预测截止时间和样本外季度。

\subsection{评价指标定义}

设第 $t$ 个季度的实际值为 $y_t$，经营桥估算为 $\widehat{y}_t$，朴素对照为 $y_t^{(0)}$，共同有效样本数为 $N$。

\paragraph{平均绝对误差（MAE）}

\begin{equation*}
MAE=\frac{1}{N}\sum_{t=1}^{N}\left|\widehat{y}_t-y_t\right|.
\end{equation*}

MAE 与目标单位相同，直接表示一次季度估算平均偏离多少，是本实验的主要指标。

\paragraph{均方根误差（RMSE）}

\begin{equation*}
RMSE=\sqrt{\frac{1}{N}\sum_{t=1}^{N}\left(\widehat{y}_t-y_t\right)^2}.
\end{equation*}

RMSE 对极端误差赋予更高权重。若 RMSE 明显高于 MAE，说明少数异常季度对结果影响较大。

\paragraph{平均偏差（Bias）}

\begin{equation*}
Bias=\frac{1}{N}\sum_{t=1}^{N}\left(\widehat{y}_t-y_t\right).
\end{equation*}

正值表示系统性高估，负值表示系统性低估。较低的 MAE 不能替代偏差检查，因为长期单向高估可能在投资判断中形成稳定误导。

\paragraph{加权绝对百分比误差（WAPE）}

\begin{equation*}
WAPE=\frac{\sum_{t=1}^{N}\left|\widehat{y}_t-y_t\right|}
{\sum_{t=1}^{N}\left|y_t\right|}.
\end{equation*}

WAPE 用于跨量级理解误差。这里不使用逐期 MAPE，因为部分产品单位毛利或公司毛利润可能接近零，逐期百分比误差会失真。

\paragraph{方向准确率}

\begin{equation*}
DirectionAccuracy=
\frac{1}{N^*}\sum_{t\in\mathcal{T}^*}
\mathbf{1}\left[
\operatorname{sgn}(y_t-y_t^{(0)})
=\operatorname{sgn}(\widehat{y}_t-y_t^{(0)})
\right],
\end{equation*}

其中 $\mathcal{T}^*$ 排除实际值与上一期完全相同的季度。该指标回答经营桥是否正确判断了“相对上一期上升还是下降”。

\paragraph{相对朴素对照改善率}

\begin{equation*}
Improvement=1-\frac{MAE_{bridge}}{MAE_{naive}}.
\end{equation*}

改善率为正表示经营桥 MAE 更低；为零表示没有改善；为负表示经营桥不如直接使用上一期实际值。例如 $-126.1\%$ 表示经营桥误差约为朴素对照的 2.261 倍，不是准确率为负。

\subsection{实验记录与后续比较契约}

每条预测由四个业务键唯一确定：\texttt{period\_end}、
\texttt{task\_group}、\texttt{target} 和 \texttt{product}。

同时锁定以下实验元数据：预测时间、训练开始季度、训练结束季度、训练行数和最后一条训练数据的公开时间。

后续候选方法可以只研究部分目标，但对于选中的每一个目标，必须：

\begin{itemize}
  \item 覆盖经营桥全部共同有效季度；
  \item 使用完全相同的预测时间和训练区间；
  \item 保证最后训练数据公开时间不晚于预测时间；
  \item 不因模型无法处理异常季度而删除样本；
  \item 在同一实际值上同时报告候选方法、经营桥和朴素对照。
\end{itemize}

如果季度覆盖或实验元数据不同，结果不进入横向比较。这一约束防止通过改变截止时间、扩大训练数据或选择容易季度制造表面改善。

\subsection{样本外结果}

滚动检验共得到 30 个样本外季度，覆盖 2018 年第四季度至 2026 年第一季度，共 450 条分任务预测。所有训练数据均早于相应预测时点，point-in-time 违规为零。

\begin{center}
\small
\begin{tabularx}{\textwidth}{L{3.1cm}Y r r r r}
\toprule
任务 & 产品 & N & 经营桥 MAE & 朴素 MAE & 改善率 \\
\midrule
实际气价 & 全产品 & 30 & 0.430 & 0.703 & 38.8\% \\
毛利润 & 全产品 & 29 & 192.923 & 229.724 & 16.0\% \\
实现售价 & 颗粒尿素 & 29 & 33.361 & 45.276 & 26.3\% \\
实现售价 & 氨 & 29 & 92.310 & 97.207 & 5.0\% \\
实现售价 & UAN & 29 & 48.476 & 43.103 & $-12.5\%$ \\
实现售价 & 硝酸铵 & 29 & 34.180 & 30.862 & $-10.8\%$ \\
单位毛利 & 颗粒尿素 & 29 & 37.887 & 44.793 & 15.4\% \\
单位毛利 & 氨 & 29 & 86.978 & 96.034 & 9.4\% \\
单位毛利 & UAN & 29 & 47.936 & 41.207 & $-16.3\%$ \\
单位毛利 & 硝酸铵 & 29 & 38.191 & 31.276 & $-22.1\%$ \\
销量 & 氨 & 30 & 135.992 & 267.467 & 49.2\% \\
销量 & 颗粒尿素 & 30 & 128.743 & 172.600 & 25.4\% \\
销量 & UAN & 30 & 154.799 & 241.300 & 35.8\% \\
销量 & 硝酸铵 & 30 & 107.232 & 47.433 & $-126.1\%$ \\
销量 & 其他产品 & 30 & 42.177 & 37.333 & $-13.0\%$ \\
\bottomrule
\end{tabularx}
\normalsize
\end{center}

实际气价的 MAE 单位为美元/MMBtu，售价和单位毛利为美元/短吨，销量为千短吨，毛利润为百万美元。正改善率表示经营桥优于上一期实际值，负值表示经营桥在该子任务上没有增加预测价值。

完整的误差和方向指标如下：

\scriptsize
\begin{longtable}{L{2.1cm}L{2.4cm}r r r r r r}
\toprule
任务 & 产品 & N & MAE & RMSE & Bias & WAPE & 方向准确率 \\
\midrule
\endfirsthead
\toprule
任务 & 产品 & N & MAE & RMSE & Bias & WAPE & 方向准确率 \\
\midrule
\endhead
实际气价 & 全产品 & 30 & 0.430 & 0.638 & 0.048 & 11.6\% & 83.3\% \\
毛利润 & 全产品 & 29 & 192.923 & 249.801 & $-7.427$ & 31.4\% & 75.9\% \\
实现售价 & 氨 & 29 & 92.310 & 120.456 & $-11.441$ & 19.5\% & 65.5\% \\
实现售价 & 硝酸铵 & 29 & 34.180 & 52.953 & $-15.524$ & 11.0\% & 72.4\% \\
实现售价 & 颗粒尿素 & 29 & 33.361 & 49.165 & $-5.452$ & 8.5\% & 79.3\% \\
实现售价 & UAN & 29 & 48.476 & 64.033 & $-19.479$ & 17.3\% & 67.9\% \\
销量 & 氨 & 30 & 135.992 & 175.151 & $-73.925$ & 14.5\% & 90.0\% \\
销量 & 硝酸铵 & 30 & 107.232 & 127.379 & 98.032 & 25.6\% & 55.2\% \\
销量 & 颗粒尿素 & 30 & 128.743 & 173.004 & $-56.535$ & 11.2\% & 80.0\% \\
销量 & 其他产品 & 30 & 42.177 & 53.299 & $-25.055$ & 7.7\% & 63.3\% \\
销量 & UAN & 30 & 154.799 & 188.272 & $-4.335$ & 9.0\% & 80.0\% \\
单位毛利 & 氨 & 29 & 86.978 & 111.664 & 16.600 & 54.6\% & 69.0\% \\
单位毛利 & 硝酸铵 & 29 & 38.191 & 54.945 & 18.901 & 60.3\% & 60.7\% \\
单位毛利 & 颗粒尿素 & 29 & 37.887 & 56.744 & 11.573 & 21.3\% & 86.2\% \\
单位毛利 & UAN & 29 & 47.936 & 60.119 & $-6.751$ & 42.7\% & 65.5\% \\
\bottomrule
\end{longtable}
\normalsize

结果表明，简单传导关系对实际气价、颗粒尿素、主要产品销量和公司毛利润具有一定样本外信息；但 UAN 和硝酸铵的售价或单位毛利尚未得到改善，硝酸铵季节销量均值明显弱于直接使用上一期销量。这些负结果不应删除，因为它们指出了基差、产品结构、披露口径或离散事件可能更重要的环节。

从 Bias 还可以看到更具体的问题：氨和颗粒尿素销量存在系统性低估，而硝酸铵销量存在系统性高估；四类产品单位毛利的 WAPE 仍在约 21\% 至 60\% 之间，说明毛利润改善并不代表每个单位经济环节都已经估计准确。

\subsection{统计解释边界}

当前结果是一组固定的样本外基准，不等同于已经证明统计显著或具备投资价值：

\begin{itemize}
  \item 每个目标只有 29---30 个样本外季度，单个异常季度可能明显改变平均误差；
  \item 同时观察 15 个子任务会产生多重比较问题，不能只挑选改善最大的指标；
  \item 扩展窗口的季度误差可能存在序列相关，普通独立样本显著性检验并不充分；
  \item 经营桥参数来自经济逻辑和前期描述性分析，仍存在研究者选择自由度；
  \item 经营预测改善不自动等于股票收益预测改善，二者需要不同的目标和执行时点。
\end{itemize}

后续候选方法除报告平均改善外，还应报告逐季度配对误差、胜率、改善是否集中于少数季度，以及适合小样本时间序列的置信区间。只有方向稳定、误差改善分布较广且经济含义合理时，才可以把结果解释为较可信的增量。

\subsection{后续方法如何证明增量价值}

后续线性回归、树模型或时间序列方法可以复用同一组实验索引。比较时至少满足：

\begin{enumerate}
  \item 使用相同的季度末后第 15 天作为信息截止点；
  \item 使用相同的扩展训练窗口和样本外季度；
  \item 对选定目标覆盖经营桥的全部有效季度，不能遗漏困难样本；
  \item 同时比较 MAE、RMSE、方向准确率和误差偏差；
  \item 先比较各中间经营变量，再比较最终毛利润，避免只看一个汇总指标；
  \item 检查改善是否集中在少数异常季度，并报告小样本不确定性。
\end{enumerate}

如果新方法只提高样本内拟合，或只在删去困难季度后优于经营桥，就不能认为它提供了增量价值。更有意义的改进方式是预测经营桥的残差，例如 UAN 地区基差、硝酸铵产品结构、检修导致的销量偏差，而不是让复杂方法重新学习已经明确的会计恒等式。

\section{领先滞后关系的统计检验}

\subsection{研究问题与两类时间时钟}

经营桥说明了变量之间应当如何连接，领先滞后检验进一步回答两个不同问题：

\begin{description}[style=nextline]
  \item[经济期时钟] 按月度或季度所属期间对齐，观察外部价格如何传导到 CF 实现售价、实际气价和盈利。它描述经营传导，但不代表市场当时已经看到 CF 财务结果；
  \item[可用期时钟] 只在来源公开时间不晚于信号日期时使用该信号，再观察之后的 CF 股票收益。它可以用于研究可执行信号，但仍不等于已经形成交易规则。
\end{description}

两类结果不能混合解释。例如“全球尿素与 CF 当季实现价同期相关”是经济关系，不意味着投资者在季度开始时已经知道 CF 当季实现价。

首版实验的分析截止日固定为 2026 年 6 月 30 日，排除尚未结束的 2026 年 7 月。共预先声明 14 组关系和 88 个滞后组合，其中 9 组使用经济期时钟，5 组使用可用期时钟。

\subsection{变量变换与滞后符号}

统一定义信号 $x_t$ 与目标 $y_{t+k}$ 配对：

\begin{equation*}
\rho_k=\operatorname{Corr}\!\left(g_x(x_t),g_y(y_{t+k})\right),
\qquad k\geq 0.
\end{equation*}

$k>0$ 表示信号领先目标；$k=0$ 表示同期关系。月频搜索 1---12 个月，季度经营关系搜索 0---4 个季度。为减少共同趋势导致的伪相关：

\begin{itemize}
  \item 尿素、天然气、产品实现价、价差和 EBITDA 主要使用一阶变化；
  \item 玉米与 CF 销量使用四季度同比变化，以降低季节性影响；
  \item 已披露毛利变化保持原本的季度环比含义，只在新财报首次进入月频面板时取一个事件样本；
  \item 股票目标为 $t+k$ 月的单月收益，不是从 $t$ 到 $t+k$ 的累计收益。
\end{itemize}

最后一条尤其重要：“领先五个月”表示信号与第五个月的单月收益相关，并不表示未来五个月累计涨幅。

\subsection{推断、稳定性和多重比较}

每个滞后组合先把信号和目标标准化，再估计一元回归；此时斜率在数值上等于相关系数。月频使用 3 阶、季度使用 2 阶 Newey--West HAC 标准误，以减轻异方差和短期序列相关对推断的影响。

同一关系会同时检查多个滞后，因此原始 $p$ 值不能直接用于挑选最佳滞后。实验同时计算关系内和全局 Benjamini--Hochberg $q$ 值。最佳滞后仅定义为预设网格中绝对相关系数最大者，仍需检查：

\begin{enumerate}
  \item 有效样本不少于月频 36 个或季度 24 个；
  \item 前半段和后半段相关系数方向一致；
  \item 两段绝对相关系数均不低于 0.10；
  \item 方向符合事前经济假设；
  \item 关系内多重检验校正后是否仍有统计证据。
\end{enumerate}

据此使用 \texttt{STRONG}、\texttt{DIRECTIONAL}、\texttt{CONTRADICTORY}、\texttt{UNSTABLE} 和 \texttt{INSUFFICIENT} 五类候选证据标签。其中 \texttt{STRONG} 要求关系内和全局校正后均达到门槛。另设 \texttt{DIAGNOSTIC} 表示变量之间存在公式包含关系，只用于核对计算方向。标签不是因果强度或交易评级。

\subsection{首版结果}

\small
\begin{longtable}{L{5.4cm}r r r r L{1.8cm}}
\toprule
关系 & 领先期 & N & 相关 & 关系内 $q$ & 标签 \\
\midrule
\endfirsthead
\toprule
关系 & 领先期 & N & 相关 & 关系内 $q$ & 标签 \\
\midrule
\endhead
全球尿素变化 $\rightarrow$ CF 氨实现价变化 & 1 季 & 40 & 0.544 & $<0.001$ & STRONG \\
全球尿素变化 $\rightarrow$ CF 颗粒尿素实现价变化 & 同期 & 40 & 0.770 & $<0.001$ & STRONG \\
全球尿素变化 $\rightarrow$ CF UAN 实现价变化 & 1 季 & 40 & 0.663 & $<0.001$ & STRONG \\
全球尿素变化 $\rightarrow$ CF 硝酸铵实现价变化 & 1 季 & 40 & 0.604 & $<0.001$ & STRONG \\
Henry Hub 变化 $\rightarrow$ CF 实际气价变化 & 1 季 & 39 & 0.640 & 0.001 & STRONG \\
全球气价差变化 $\rightarrow$ CF 实现气价差变化 & 同期 & 39 & 0.775 & $<0.001$ & STRONG \\
全球气价差变化 $\rightarrow$ CF 毛利率变化 & 同期 & 41 & 0.366 & $<0.001$ & STRONG \\
全球气价差变化 $\rightarrow$ CF EBITDA 变化 & 同期 & 41 & 0.633 & $<0.001$ & STRONG \\
玉米同比变化 $\rightarrow$ CF 销量同比变化 & 1 季 & 37 & $-0.375$ & 0.087 & CONTRADICTORY \\
Henry Hub 变化 $\rightarrow$ 全球气价差变化 & 同期 & 157 & $-0.416$ & 0.001 & DIAGNOSTIC \\
已披露毛利变化 $\rightarrow$ CF 单月收益 & 5 月 & 39 & 0.325 & 0.005 & STRONG \\
已披露实现气价差变化 $\rightarrow$ CF 单月收益 & 5 月 & 37 & 0.360 & 0.001 & UNSTABLE \\
全球气价差变化 $\rightarrow$ CF 单月收益 & 8 月 & 153 & 0.141 & 0.076 & UNSTABLE \\
玉米变化 $\rightarrow$ CF 单月收益 & 12 月 & 149 & $-0.118$ & 0.838 & UNSTABLE \\
\bottomrule
\end{longtable}
\normalsize

\begin{figure}[htbp]
\centering
% Generated by scripts/generate_lead_lag_figure.py; do not edit manually.
% Source analysis: cf_lead_lag_v1 v1.0.0
\definecolor{cfStrong}{HTML}{0072B2}
\definecolor{cfDirectional}{HTML}{009E73}
\definecolor{cfContradictory}{HTML}{D55E00}
\definecolor{cfUnstable}{HTML}{999999}
\definecolor{cfInsufficient}{HTML}{CCCCCC}
\definecolor{cfDiagnostic}{HTML}{7B3294}
\begin{tikzpicture}
\begin{axis}[
  width=0.63\textwidth,
  height=15.2cm,
  xmin=-0.80, xmax=1.08,
  ymin=0.35, ymax=14.85,
  axis x line*=bottom,
  axis y line*=left,
  y axis line style={draw=none},
  ytick style={draw=none},
  xlabel={最佳滞后的相关系数 / 标准化回归斜率},
  xtick={-0.75,-0.50,-0.25,0,0.25,0.50,0.75,1.00},
  grid=both,
  major grid style={black!12},
  minor grid style={black!5},
  ytick={1,2,3,4,5,6,7,8,9,10,11,12,13,14},
  yticklabels={{Henry Hub $\rightarrow$ 全球气价差（同期）},{玉米 $\rightarrow$ CF 单月收益（12 月）},{全球气价差 $\rightarrow$ CF 单月收益（8 月）},{已披露实现气价差 $\rightarrow$ CF 单月收益（5 月）},{已披露毛利变化 $\rightarrow$ CF 单月收益（5 月）},{玉米 $\rightarrow$ CF 销量同比（1 季）},{全球气价差 $\rightarrow$ CF EBITDA 变化（同期）},{全球气价差 $\rightarrow$ CF 毛利率变化（同期）},{全球气价差 $\rightarrow$ CF 实现气价差（同期）},{Henry Hub $\rightarrow$ CF 实际气价（1 季）},{全球尿素 $\rightarrow$ 硝酸铵实现价（1 季）},{全球尿素 $\rightarrow$ UAN 实现价（1 季）},{全球尿素 $\rightarrow$ 氨实现价（1 季）},{全球尿素 $\rightarrow$ 颗粒尿素实现价（同期）}},
  yticklabel style={font=\scriptsize,align=right,text width=6.0cm},
  tick label style={font=\scriptsize},
  label style={font=\small},
  legend style={at={(0.5,-0.10)},anchor=north,draw=none,font=\scriptsize,/tikz/every even column/.append style={column sep=0.18cm}},
  legend columns=4,
]
\addplot[black!55,thin,forget plot] coordinates {(0,0.5) (0,14.65)};
\addplot[black!25,dashed,forget plot] coordinates {(-0.80,5.5) (1.08,5.5)};
\node[anchor=west,font=\scriptsize\bfseries,text=black!70] at (axis cs:-0.78,14.55) {经济期经营传导};
\node[anchor=west,font=\scriptsize\bfseries,text=black!70] at (axis cs:-0.78,5.35) {可用期与市场反应};
\addplot[
  only marks,
  color=cfStrong, mark=*, mark size=2.4pt,
  error bars/.cd, x dir=both, x explicit,
  error bar style={line width=0.7pt},
  error mark options={rotate=90,mark size=2pt,line width=0.7pt},
] coordinates {
    (0.769565,14) +- (0.150067,0)
    (0.543812,13) +- (0.168423,0)
    (0.662938,12) +- (0.162848,0)
    (0.603850,11) +- (0.272433,0)
    (0.640412,10) +- (0.363364,0)
    (0.774650,9) +- (0.222231,0)
    (0.366085,8) +- (0.144392,0)
    (0.632522,7) +- (0.218913,0)
    (0.325344,5) +- (0.208880,0)
};
\addlegendentry{STRONG}
\addplot[
  only marks,
  color=cfContradictory, mark=diamond*, mark size=2.4pt,
  error bars/.cd, x dir=both, x explicit,
  error bar style={line width=0.7pt},
  error mark options={rotate=90,mark size=2pt,line width=0.7pt},
] coordinates {
    (-0.374834,6) +- (0.335234,0)
};
\addlegendentry{CONTRADICTORY}
\addplot[
  only marks,
  color=cfUnstable, mark=*, mark size=2.4pt,
  error bars/.cd, x dir=both, x explicit,
  error bar style={line width=0.7pt},
  error mark options={rotate=90,mark size=2pt,line width=0.7pt},
] coordinates {
    (0.359837,4) +- (0.184396,0)
    (0.141320,3) +- (0.101512,0)
    (-0.117814,2) +- (0.137475,0)
};
\addlegendentry{UNSTABLE}
\addplot[
  only marks,
  color=cfDiagnostic, mark=square*, mark size=2.4pt,
  error bars/.cd, x dir=both, x explicit,
  error bar style={line width=0.7pt},
  error mark options={rotate=90,mark size=2pt,line width=0.7pt},
] coordinates {
    (-0.415567,1) +- (0.207785,0)
};
\addlegendentry{DIAGNOSTIC}
\end{axis}
\end{tikzpicture}

\caption{M4.1 各关系最佳滞后的标准化系数与 HAC 95\% 区间。括号中的月份或季度表示信号领先目标的期间；股票收益对应指定月份的单月收益，不是累计收益。颜色表示完整报告中的证据标签。由于每个点是在预设滞后网格中按绝对相关系数选出的最佳结果，区间本身没有进一步校正“最佳滞后选择”的不确定性。}
\label{fig:cf-lead-lag-best-coefficients}
\end{figure}

经营端的结果与产业逻辑大体一致：颗粒尿素实现价更接近全球尿素同期变化，氨、UAN 和硝酸铵则表现出约一个季度的传导滞后；CF 实际气价对 Henry Hub 也有约一季度滞后；利润代理向实现价差和 EBITDA 的关系主要发生在同期。

Henry Hub 与全球尿素气价差的同期负相关不算独立证据，因为该价差的计算公式本身就扣除了 Henry Hub 天然气成本。保留这一行的作用是验证指标构造方向，不能把它再次计入周期状态评分。

需求端不能从玉米价格直接推出 CF 销量。当前玉米与销量的最佳相关方向为负，虽然前后子样本方向一致，但不符合事前假设且校正后证据不足，不能放入核心状态机。可能原因包括 CF 销量受产能、检修、库存、出口和季节调度约束，而玉米价格只是农户经济性的一个维度。

股票端更需要克制。只有“最新披露毛利变化与第五个月单月收益”通过首轮门槛，但事件样本仅 39 个，而且最佳月份是在六个候选滞后中选择出来的。外部气价差、已披露实现价差和玉米到股票收益的关系没有通过前后样本稳定性检查。在完成下一步滚动窗口、季节和压力分组前，不把这些关系直接写入买入评分。

\subsection{固定滞后的稳定性检验}

M4.2 不再寻找新的最佳滞后，而是把 M4.1 已选出的滞后固定下来，再回答“这条关系是否只在某一段历史或某一种场景成立”。这样可以避免在每个子样本中重新挑选最有利的滞后。首版包含：

\begin{itemize}
  \item 月频连续关系使用 60 个观察的滚动窗口，季度关系使用 20 个观察，财报事件关系使用 20 个事件；
  \item 以目标期间划分春季施肥、秋季施肥和淡季；季度样本只分为上半年与下半年，避免每季只有约 10 个观察；
  \item 以对齐样本的 Henry Hub 中位数划分高低气价；
  \item 以全球尿素气价差绝对变化的 75 分位数划分正常变化和大幅变化；
  \item 滚动窗口预期方向占比至少为 75\%，滚动中位绝对相关至少为 0.20；
  \item 高低气价和正常/大幅变化场景的预期方向占比至少为 75\%，每个场景绝对相关至少为 0.10。
\end{itemize}

季节切片不作为硬门槛，因为季度样本分组后只有约 19---20 个观察。滚动窗口彼此高度重叠，因此窗口同向占比是稳定性描述，不是独立试验的成功概率。高低气价和大幅变化阈值使用全样本分位数，也只用于历史稳健性检验，不能直接作为实时状态阈值。

\subsection{M4.2 筛选结果}

固定滞后后共得到 506 个滚动窗口和 89 个场景切片，重复运行结果一致。关系按用途而不是单纯按相关系数分类：

\begin{center}
\small
\begin{tabularx}{\textwidth}{L{3.4cm}r Y}
\toprule
决策 & 数量 & 含义和关系 \\
\midrule
\texttt{ACCEPT\_BRIDGE} & 5 & 四类全球尿素到 CF 产品实现价，以及 Henry Hub 到 CF 实际气价；保留在经营桥中 \\
\texttt{ACCEPT\_CORE\_VALIDATION} & 3 & 全球尿素气价差到 CF 实现价差、毛利率变化和 EBITDA 变化；共同验证一个核心代理，不重复计三次权重 \\
\texttt{ACCEPT\_CONFIRMATION} & 1 & 最新披露毛利变化到第五个月单月收益；只作为市场确认 \\
\texttt{REJECT} & 4 & 玉米到销量、玉米到股票收益、外部气价差到股票收益、已披露实现价差到股票收益 \\
\texttt{DIAGNOSTIC} & 1 & Henry Hub 到理论气价差的公式包含关系 \\
\bottomrule
\end{tabularx}
\normalsize
\end{center}

\begin{figure}[htbp]
\centering
% Generated by scripts/generate_stability_figure.py; do not edit manually.
% Source analysis: cf_lead_lag_stability_v1 v1.0.0
\definecolor{cfBridge}{HTML}{0072B2}
\definecolor{cfCore}{HTML}{005A8D}
\definecolor{cfConfirmation}{HTML}{009E73}
\definecolor{cfDemand}{HTML}{56B4E9}
\definecolor{cfConditional}{HTML}{E69F00}
\definecolor{cfReject}{HTML}{D55E00}
\definecolor{cfDiagnostic}{HTML}{7B3294}
\begin{tikzpicture}
\begin{axis}[
  width=0.63\textwidth,
  height=15.2cm,
  xmin=-0.80, xmax=1.08,
  ymin=0.35, ymax=14.85,
  axis x line*=bottom,
  axis y line*=left,
  y axis line style={draw=none},
  ytick style={draw=none},
  xlabel={固定滞后下的相关系数},
  xtick={-0.75,-0.50,-0.25,0,0.25,0.50,0.75,1.00},
  grid=both,
  major grid style={black!12},
  minor grid style={black!5},
  ytick={1,2,3,4,5,6,7,8,9,10,11,12,13,14},
  yticklabels={{Henry Hub $\rightarrow$ 全球气价差（同期）},{玉米 $\rightarrow$ CF 单月收益（12 月）},{全球气价差 $\rightarrow$ CF 单月收益（8 月）},{已披露实现气价差 $\rightarrow$ CF 单月收益（5 月）},{已披露毛利变化 $\rightarrow$ CF 单月收益（5 月）},{玉米 $\rightarrow$ CF 销量同比（1 季）},{全球气价差 $\rightarrow$ CF EBITDA 变化（同期）},{全球气价差 $\rightarrow$ CF 毛利率变化（同期）},{全球气价差 $\rightarrow$ CF 实现气价差（同期）},{Henry Hub $\rightarrow$ CF 实际气价（1 季）},{全球尿素 $\rightarrow$ 硝酸铵实现价（1 季）},{全球尿素 $\rightarrow$ UAN 实现价（1 季）},{全球尿素 $\rightarrow$ 氨实现价（1 季）},{全球尿素 $\rightarrow$ 颗粒尿素实现价（同期）}},
  yticklabel style={font=\scriptsize,align=right,text width=6.0cm},
  tick label style={font=\scriptsize},
  label style={font=\small},
  legend style={at={(0.5,-0.10)},anchor=north,draw=none,font=\scriptsize,/tikz/every even column/.append style={column sep=0.16cm}},
  legend columns=3,
]
\addplot[black!55,thin,forget plot] coordinates {(0,0.5) (0,14.65)};
\addplot[black!25,dashed,forget plot] coordinates {(-0.80,5.5) (1.08,5.5)};
\node[anchor=west,font=\scriptsize\bfseries,text=black!70] at (axis cs:-0.78,14.55) {经济期经营传导};
\node[anchor=west,font=\scriptsize\bfseries,text=black!70] at (axis cs:-0.78,5.35) {可用期与市场反应};
\addplot[cfBridge,line width=1.0pt,forget plot] coordinates {
  (0.620933,14) (0.882138,14)
};
\addplot[cfBridge,line width=1.0pt,forget plot] coordinates {
  (0.226974,13) (0.658547,13)
};
\addplot[cfBridge,line width=1.0pt,forget plot] coordinates {
  (0.593192,12) (0.851506,12)
};
\addplot[cfBridge,line width=1.0pt,forget plot] coordinates {
  (0.600383,11) (0.800729,11)
};
\addplot[cfBridge,line width=1.0pt,forget plot] coordinates {
  (0.437936,10) (0.699459,10)
};
\addplot[cfCore,line width=1.0pt,forget plot] coordinates {
  (0.699113,9) (0.870230,9)
};
\addplot[cfCore,line width=1.0pt,forget plot] coordinates {
  (0.278461,8) (0.670505,8)
};
\addplot[cfCore,line width=1.0pt,forget plot] coordinates {
  (0.060195,7) (0.754394,7)
};
\addplot[cfReject,line width=1.0pt,forget plot] coordinates {
  (-0.668020,6) (-0.479311,6)
};
\addplot[cfConfirmation,line width=1.0pt,forget plot] coordinates {
  (-0.051343,5) (0.543231,5)
};
\addplot[cfReject,line width=1.0pt,forget plot] coordinates {
  (-0.056306,4) (0.518208,4)
};
\addplot[cfReject,line width=1.0pt,forget plot] coordinates {
  (0.021112,3) (0.262281,3)
};
\addplot[cfReject,line width=1.0pt,forget plot] coordinates {
  (-0.280556,2) (-0.016528,2)
};
\addplot[cfDiagnostic,line width=1.0pt,forget plot] coordinates {
  (-0.732878,1) (-0.343332,1)
};
\addplot[only marks,color=cfBridge,mark=*,mark size=2.5pt] coordinates {
  (0.837398,14)
  (0.566758,13)
  (0.675084,12)
  (0.722427,11)
  (0.641723,10)
};
\addlegendentry{滚动中位：经营桥}
\addplot[only marks,color=cfCore,mark=square*,mark size=2.5pt] coordinates {
  (0.821152,9)
  (0.575228,8)
  (0.714412,7)
};
\addlegendentry{滚动中位：核心验证}
\addplot[only marks,color=cfConfirmation,mark=triangle*,mark size=2.5pt] coordinates {
  (0.236424,5)
};
\addlegendentry{滚动中位：市场确认}
\addplot[only marks,color=cfReject,mark=x,mark size=2.5pt] coordinates {
  (-0.551449,6)
  (0.341482,4)
  (0.169237,3)
  (-0.176157,2)
};
\addlegendentry{滚动中位：拒绝}
\addplot[only marks,color=cfDiagnostic,mark=square*,mark size=2.5pt] coordinates {
  (-0.411920,1)
};
\addlegendentry{滚动中位：诊断}
\addplot[only marks,color=black,mark=diamond,mark size=2.8pt,mark options={fill=white}] coordinates {
  (0.769565,14)
  (0.543812,13)
  (0.662938,12)
  (0.603850,11)
  (0.640412,10)
  (0.774650,9)
  (0.366085,8)
  (0.632522,7)
  (-0.374834,6)
  (0.325344,5)
  (0.359837,4)
  (0.141320,3)
  (-0.117814,2)
  (-0.415567,1)
};
\addlegendentry{全样本相关}
\end{axis}
\end{tikzpicture}

\caption{M4.2 固定滞后下的滚动相关范围。彩色横线表示所有滚动窗口相关系数的最小值至最大值，彩色标记表示滚动中位数，空心菱形表示全样本相关。窗口高度重叠，范围只用于观察参数稳定性，不构成额外独立样本。}
\label{fig:cf-lead-lag-rolling-stability}
\end{figure}

经营桥关系的滚动方向均为 100\% 一致。颗粒尿素实现价传导的滚动中位相关为 0.837；Henry Hub 到 CF 实际气价为 0.642；四类售价关系在高低气价和正常/大幅价差变化中均保持预期方向。因此，它们可以进入经营估算，但仍按产品权重汇总，不能作为五个独立周期信号重复计分。

全球尿素气价差到 CF 实现价差、毛利率变化和 EBITDA 变化也全部通过。三者的滚动中位相关分别约为 0.821、0.575 和 0.714。这意味着全球尿素气价差可以作为核心经营周期代理；三条下游关系是对该代理的验证，而不是三项新增特征。

场景强度仍有明显差异。Henry Hub 到 CF 实际气价在低气价样本中的相关约为 0.952，在高气价样本中降至约 0.372；颗粒尿素售价传导在大幅价差变化时约为 0.888，在正常变化时约为 0.371。方向没有翻转，但线性敏感度并非常数，后续状态机应把压力场景作为解释标签，而不是假定固定斜率。

股票关系继续单独处理。最新披露毛利变化的滚动窗口有 95\% 保持正方向，高低气价和大幅变化场景也同向，因此保留为市场确认；但春季施肥月份相关只有约 0.084，且目标是第五个月的单月收益，所以它不能进入经营周期核心分数。其余三条股票关系延续 M4.1 的拒绝结论。

这里的“大幅变化”只表示气价差绝对变化位于历史前 25\%，并不等于已经识别出战争、制裁、装置停产或贸易政策等具体供给事件。真正的事件窗口仍需要独立、可审计的事件日期表。

\subsection{M4.3 规则型经营周期状态}

M4.3 将 M4.2 的筛选结论固化为可逐月复算的状态机。首版刻意把经营周期和股票确认分成两层：经营状态只使用全球尿素气价差，CF 最新披露毛利变化和六个月股价动量只用于描述确认或背离，不参与状态判定。

核心价差虽然由全球尿素和 Henry Hub 两项构成，但经济上只代表一个“产品价格减主要原料成本”的利润代理。M4.2 中通过检验的 CF 实现价差、毛利率变化和 EBITDA 变化是对这一个代理的下游验证，不能再分别获得权重。否则同一条价格传导链会被重复计分。

对每个月的核心价差 $S_t$，使用截至当月的 60 个月窗口计算位置标准分，并计算三个月变化的标准分：

\begin{align*}
L_t &= z_{t,60}(S_t), \\
D_t &= z_{t,60}(S_t-S_{t-3}), \\
C_t &= \frac{1}{2}\operatorname{clip}(L_t,-2,2)
     + \frac{1}{2}\operatorname{clip}(D_t,-2,2).
\end{align*}

最低需要 36 个月历史。$C_t$ 只用于连续展示，离散状态由 $L_t$ 和 $D_t$ 的二维阈值直接确定：

\begin{center}
\small
\begin{tabularx}{\textwidth}{L{2.8cm}L{4.5cm}Y}
\toprule
状态 & 进入条件 & 经济含义 \\
\midrule
\texttt{RECOVERY} & $L_t\leq0.25,\ D_t\geq0.25$ & 价差位置尚不高，但方向已经改善 \\
\texttt{EXPANSION} & $L_t\geq0.25,\ D_t\geq0.25$ & 价差位置和方向同时偏强 \\
\texttt{PEAK\_RISK} & $L_t\geq0.75,\ D_t\leq-0.25$ & 利润代理仍高，但方向已经转弱 \\
\texttt{CONTRACTION} & $L_t\leq-0.25,\ D_t\leq-0.25$ & 价差位置和方向同时偏弱 \\
\texttt{TROUGH} & $L_t\leq-0.75,\ D_t\geq0.25$ & 价差仍低，但方向开始修复 \\
\texttt{MIXED} & 其余组合或核心数据不可用 & 证据冲突、历史不足或数据缺口 \\
\bottomrule
\end{tabularx}
\normalsize
\end{center}

新状态需要连续两个月满足条件；非 \texttt{MIXED} 状态原则上至少持续三个月，并用较宽的退出条件形成滞回。唯一的安全例外是核心数据缺失：缺失当月立即进入 \texttt{MIXED}，不能用最短持续期保留旧状态。

\begin{figure}[htbp]
\centering
\definecolor{stateRecovery}{HTML}{8DD3C7}
\definecolor{stateExpansion}{HTML}{80B1D3}
\definecolor{statePeak}{HTML}{FDB462}
\definecolor{stateContraction}{HTML}{FB8072}
\definecolor{stateTrough}{HTML}{B3DE69}
\definecolor{stateMixed}{HTML}{D9D9D9}
\begin{tikzpicture}
\begin{axis}[
  width=\textwidth,height=7.1cm,
  xmin=2016,xmax=2027,ymin=-2.1,ymax=2.1,
  xtick={2016,2018,2020,2022,2024,2026},
  ytick={-2,-1,0,1,2},
  xlabel={月份},ylabel={周期分数},
  axis lines=left,grid=major,grid style={gray!18},
  tick label style={font=\small},label style={font=\small},
  clip=false
]
\path[fill=stateMixed,fill opacity=0.32,draw=none] (axis cs:2016.00000,-2.1) rectangle (axis cs:2016.25000,2.1);
\path[fill=stateContraction,fill opacity=0.32,draw=none] (axis cs:2016.25000,-2.1) rectangle (axis cs:2016.83333,2.1);
\path[fill=stateTrough,fill opacity=0.32,draw=none] (axis cs:2016.83333,-2.1) rectangle (axis cs:2017.41667,2.1);
\path[fill=stateContraction,fill opacity=0.32,draw=none] (axis cs:2017.41667,-2.1) rectangle (axis cs:2017.75000,2.1);
\path[fill=stateTrough,fill opacity=0.32,draw=none] (axis cs:2017.75000,-2.1) rectangle (axis cs:2018.33333,2.1);
\path[fill=stateRecovery,fill opacity=0.32,draw=none] (axis cs:2018.33333,-2.1) rectangle (axis cs:2019.00000,2.1);
\path[fill=stateExpansion,fill opacity=0.32,draw=none] (axis cs:2019.00000,-2.1) rectangle (axis cs:2019.25000,2.1);
\path[fill=stateMixed,fill opacity=0.32,draw=none] (axis cs:2019.25000,-2.1) rectangle (axis cs:2019.66667,2.1);
\path[fill=stateExpansion,fill opacity=0.32,draw=none] (axis cs:2019.66667,-2.1) rectangle (axis cs:2019.91667,2.1);
\path[fill=stateMixed,fill opacity=0.32,draw=none] (axis cs:2019.91667,-2.1) rectangle (axis cs:2020.33333,2.1);
\path[fill=stateExpansion,fill opacity=0.32,draw=none] (axis cs:2020.33333,-2.1) rectangle (axis cs:2020.58333,2.1);
\path[fill=stateContraction,fill opacity=0.32,draw=none] (axis cs:2020.58333,-2.1) rectangle (axis cs:2020.83333,2.1);
\path[fill=stateExpansion,fill opacity=0.32,draw=none] (axis cs:2020.83333,-2.1) rectangle (axis cs:2022.25000,2.1);
\path[fill=statePeak,fill opacity=0.32,draw=none] (axis cs:2022.25000,-2.1) rectangle (axis cs:2022.91667,2.1);
\path[fill=stateExpansion,fill opacity=0.32,draw=none] (axis cs:2022.91667,-2.1) rectangle (axis cs:2023.16667,2.1);
\path[fill=stateMixed,fill opacity=0.32,draw=none] (axis cs:2023.16667,-2.1) rectangle (axis cs:2023.83333,2.1);
\path[fill=stateRecovery,fill opacity=0.32,draw=none] (axis cs:2023.83333,-2.1) rectangle (axis cs:2024.08333,2.1);
\path[fill=stateMixed,fill opacity=0.32,draw=none] (axis cs:2024.08333,-2.1) rectangle (axis cs:2024.41667,2.1);
\path[fill=stateContraction,fill opacity=0.32,draw=none] (axis cs:2024.41667,-2.1) rectangle (axis cs:2024.83333,2.1);
\path[fill=stateMixed,fill opacity=0.32,draw=none] (axis cs:2024.83333,-2.1) rectangle (axis cs:2025.08333,2.1);
\path[fill=stateContraction,fill opacity=0.32,draw=none] (axis cs:2025.08333,-2.1) rectangle (axis cs:2025.33333,2.1);
\path[fill=stateRecovery,fill opacity=0.32,draw=none] (axis cs:2025.33333,-2.1) rectangle (axis cs:2026.00000,2.1);
\path[fill=stateContraction,fill opacity=0.32,draw=none] (axis cs:2026.00000,-2.1) rectangle (axis cs:2026.16667,2.1);
\path[fill=stateMixed,fill opacity=0.32,draw=none] (axis cs:2026.16667,-2.1) rectangle (axis cs:2026.50000,2.1);
\addplot[black,very thick] coordinates {
  (2016.07527,nan)
  (2016.15860,nan)
  (2016.24731,-1.443533)
  (2016.32527,-1.789656)
  (2016.41398,-1.507221)
  (2016.49462,-1.279891)
  (2016.57527,-1.626944)
  (2016.66398,-1.611737)
  (2016.74462,-1.139051)
  (2016.83065,-0.404270)
  (2016.91129,-0.139415)
  (2016.99462,-0.652421)
  (2017.08065,-0.163897)
  (2017.15591,0.239219)
  (2017.24731,-0.010395)
  (2017.32258,-0.456964)
  (2017.41398,-0.695675)
  (2017.49462,-0.897675)
  (2017.58065,-0.913949)
  (2017.66398,-0.949372)
  (2017.74194,-0.559816)
  (2017.83065,0.097026)
  (2017.91129,0.753446)
  (2017.99194,1.064336)
  (2018.08065,-0.139311)
  (2018.15591,-0.471758)
  (2018.24194,-0.726935)
  (2018.32796,0.182133)
  (2018.41398,0.110264)
  (2018.49194,-0.295792)
  (2018.58065,-0.339000)
  (2018.66398,0.113734)
  (2018.73925,0.671848)
  (2018.83065,0.715732)
  (2018.91129,0.108671)
  (2018.99731,0.495219)
  (2019.08065,0.901553)
  (2019.15591,0.780908)
  (2019.24194,-0.100174)
  (2019.32796,-0.322209)
  (2019.41398,-0.170528)
  (2019.48925,0.258415)
  (2019.58065,0.258989)
  (2019.66129,0.663415)
  (2019.74462,0.570923)
  (2019.83065,0.255910)
  (2019.90860,-0.506984)
  (2019.99731,-0.412042)
  (2020.08065,-0.356985)
  (2020.15591,-0.058609)
  (2020.24731,-0.031345)
  (2020.32796,0.446452)
  (2020.40860,0.753732)
  (2020.49462,0.205901)
  (2020.58065,-0.546473)
  (2020.66398,-0.807066)
  (2020.74462,0.676614)
  (2020.82796,1.192546)
  (2020.91129,0.956156)
  (2020.99731,0.269932)
  (2021.07527,0.097174)
  (2021.15054,-1.309754)
  (2021.24731,2.000000)
  (2021.32796,2.000000)
  (2021.40591,2.000000)
  (2021.49462,1.058313)
  (2021.57796,0.875915)
  (2021.66398,1.758482)
  (2021.74462,1.964603)
  (2021.82527,1.225205)
  (2021.91129,2.000000)
  (2021.99731,2.000000)
  (2022.08065,2.000000)
  (2022.15591,2.000000)
  (2022.24731,0.548569)
  (2022.32527,0.192481)
  (2022.41398,0.678613)
  (2022.49462,0.624456)
  (2022.57527,-0.030622)
  (2022.66398,-0.510190)
  (2022.74462,-0.592549)
  (2022.83065,0.448484)
  (2022.91129,1.007939)
  (2022.99462,0.725642)
  (2023.08065,0.203119)
  (2023.15591,-0.242079)
  (2023.24731,-0.622211)
  (2023.32258,-0.933801)
  (2023.41398,-0.837228)
  (2023.49462,-0.407047)
  (2023.58065,-0.358644)
  (2023.66398,-0.205341)
  (2023.74194,0.088040)
  (2023.83065,0.260901)
  (2023.91129,0.344024)
  (2023.99194,0.158685)
  (2024.08065,-0.200706)
  (2024.15860,-0.230166)
  (2024.23925,-0.210408)
  (2024.32796,-0.081780)
  (2024.41398,-0.300888)
  (2024.48925,-0.539618)
  (2024.58065,-0.336690)
  (2024.66129,-0.059846)
  (2024.74462,0.042259)
  (2024.83065,-0.088536)
  (2024.90860,-0.039567)
  (2024.99731,-0.154091)
  (2025.08065,-0.392351)
  (2025.15591,-0.380184)
  (2025.24731,0.009875)
  (2025.32796,0.236295)
  (2025.41129,0.051690)
  (2025.49462,-0.132248)
  (2025.58065,-0.082583)
  (2025.65860,0.400533)
  (2025.74462,0.656396)
  (2025.83065,0.368794)
  (2025.90591,-0.372322)
  (2025.99731,-0.769554)
  (2026.07796,-1.156666)
  (2026.15323,-0.408600)
  (2026.24731,nan)
  (2026.32796,nan)
  (2026.40860,nan)
  (2026.49462,nan)
};
\addplot[black!45,densely dashed] coordinates {(2016,0) (2027,0)};
\end{axis}
\node[anchor=north,font=\scriptsize,yshift=-3mm] at (current bounding box.south) {
\textcolor{stateRecovery}{\rule{0.24cm}{0.15cm}} 修复\quad
\textcolor{stateExpansion}{\rule{0.24cm}{0.15cm}} 扩张\quad
\textcolor{statePeak}{\rule{0.24cm}{0.15cm}} 高位转弱\quad
\textcolor{stateContraction}{\rule{0.24cm}{0.15cm}} 收缩\quad
\textcolor{stateTrough}{\rule{0.24cm}{0.15cm}} 低位转强\quad
\textcolor{stateMixed}{\rule{0.24cm}{0.15cm}} 混合/缺口};
\end{tikzpicture}

\caption{M4.3 CF 经营周期状态时间轴。黑线为价差位置与三个月动量的等权连续分数，背景色为经过连续两个月确认、最短持续期和滞回处理后的离散状态。灰色既可能表示信号冲突，也可能表示核心数据缺失，必须结合状态原因字段区分。}
\label{fig:cf-cycle-state-timeline}
\end{figure}

历史结果给出了若干可以复盘的阶段。2016 年 4---10 月被识别为收缩，随后进入低位转强；2020 年 11 月至 2022 年 3 月为持续扩张；2022 年 4---11 月价差仍处高位但动量转负，因此进入 \texttt{PEAK\_RISK}。这一定义不使用估值，避免把 M5 的结论提前写入 M4。2018 年 5---12 月和 2025 年 5---12 月被标记为修复，但后者年末原始信号已经转弱，经过确认规则后在 2026 年 1 月进入收缩，说明状态标签相对于月度噪声有意保留了迟滞。

截至 2026 年 6 月，核心价差从 3 月起不可用，因为 World Bank 全球尿素数据停留在 2025 年 12 月，超过数据面板的 75 天有效期。当前经营周期状态因此是 \texttt{MIXED/DATA\_GAP}，而不是沿用 2 月的收缩状态。最新 CF 披露毛利率环比变化为负，CF 六个月动量为正，确认层标记为 \texttt{DIVERGENT}。这个背离可以提醒研究者继续核对新数据，但不能替代缺失的核心信号。

该状态机是后续评估的基准，不是买入评分。阈值在观察股票收益前声明，没有使用未来收益选择切点；\texttt{PEAK\_RISK} 和 \texttt{TROUGH} 也都不包含估值条件。M5 可以在这些经营状态之上研究估值，M6 再决定状态、估值和确认层如何形成观察规则。

\section{时间字段应该如何理解}

研究数据中常见三类时间：

\begin{description}[style=nextline]
  \item[观察日期] 数据描述的经济日期，例如天然气交易日、肥料报告日或财务季度末；
  \item[公开日期] 市场实际能够看到该数据的日期，例如报告发布日期或 SEC filing date；
  \item[本地获取日期] 数据被下载和保存的日期。
\end{description}

例如，一季度财务数据描述的是 3 月 31 日结束的季度，但可能在 5 月才公开。研究 4 月份的市场状态时不能使用尚未公开的一季度结果。

低频数据出现在日度表中时，表示“截至当天最新已经公开的值”。它不意味着该指标当天重新观测了一次。

\section{数据覆盖与当前限制}

\begin{center}
\small
\begin{tabularx}{\textwidth}{L{3.1cm}L{3.2cm}Y}
\toprule
数据 & 当前覆盖 & 主要限制 \\
\midrule
CF 股票 & 2013-01 至 2026-07 & 只有约 13.5 年，完整周期数量有限 \\
Henry Hub & 1997-01 至 2026-07 & 现货价不等于 CF 实际采购价 \\
全球尿素 & 1960-01 至 2025-12 & 月频且当前陈旧，地区基差明显 \\
AMS 肥料 & 2022-09 至 2026-07 & 历史短，Illinois 报价不能代表全部美国市场 \\
CF 分产品经营 & 2015Q2 至 2026Q1 & 只有约 42 个季度，披露口径可能变化 \\
农产品期货 & 2013-01 至 2026-07 & 连续合约换月会引入价格跳变 \\
种植面积 & 2013 至 2026 & 低频、多次修订，部分历史公开日期近似 \\
农业成本回报 & 1996/1997 至 2025 & 年度数据，多项目行数不等于独立样本数 \\
\bottomrule
\end{tabularx}
\normalsize
\end{center}

当前最重要的三项限制是：

\begin{enumerate}
  \item \textbf{全球尿素近期缺口：}最新 World Bank 尿素数据停留在 2025-12，近期全球价差不能继续依赖旧值；
  \item \textbf{公司季度样本少：}分产品经营数据只有约 42 个季度，适合解释和校准，不支持高自由度模型；
  \item \textbf{不同价格口径不可混同：}全球基准、Illinois 经销商报价和 CF 公司实现价分别代表不同地区和交易环节。
\end{enumerate}

\section{按研究问题选择数据}

\begin{center}
\small
\begin{tabularx}{\textwidth}{L{4.1cm}Y}
\toprule
研究问题 & 建议优先观察 \\
\midrule
全球氮肥周期处于什么位置？ & 全球尿素、Henry Hub、全球尿素气价差 \\
美国现货是否确认全球方向？ & AMS 氨/尿素/UAN、美国氮肥价格篮子 \\
农业需求是否有支撑？ & 玉米价格、玉米相对大豆价格、种植面积、ERS 农户利润 \\
外部价格是否传导到 CF？ & CF 分产品实现售价、实际气价、实现气价差、毛利/吨 \\
本季度盈利为何变化？ & 实现售价、销量、实际气价、产品结构、其他成本残差 \\
股价是否已经反映周期？ & CF 六个月动量、估值、外部价差和最新公司盈利方向 \\
春耕或秋肥是否构成短期催化？ & 施肥季节、AMS 报价、玉米价格和种植面积 \\
\bottomrule
\end{tabularx}
\normalsize
\end{center}

\section{不应从数据中直接得出的结论}

\begin{itemize}
  \item 尿素价格上涨不必然意味着 CF 毛利同比上涨；天然气、销量、库存和基差可能抵消影响；
  \item Henry Hub 下跌不必然立即改善当季利润；采购合同和库存时点会产生滞后；
  \item 玉米价格上涨不必然带来同等比例的氮肥需求增长；面积、天气和单位施肥量同样重要；
  \item 理论现金价差不是会计毛利，也不是 EBITDA；
  \item 日度表中的年度或季度字段重复出现，不代表获得了更多独立观测；
  \item 历史相关性不能替代对产能、贸易、制裁、天然气和天气冲击的情景判断。
\end{itemize}

\section{术语表}

\begin{description}[style=nextline]
  \item[氨（Ammonia）] 基础氮肥产品，也是尿素和 UAN 等下游氮肥的原料；
  \item[尿素（Urea）] 高含氮固体肥料，国际贸易活跃，是常见全球氮肥基准；
  \item[UAN] 尿素硝铵溶液，常见含氮浓度为 28\% 或 32\%；
  \item[MMBtu] 百万英热单位，天然气价格和耗用量的常见单位；
  \item[短吨] 美国常用重量单位，1 短吨为 2,000 磅；
  \item[产品吨] 按产品总重量计量，不考虑其中含氮比例；
  \item[含氮吨] 按产品中的氮元素重量计量，便于跨产品比较；
  \item[实现售价] 公司实际确认的平均销售价格，与公开现货报价不同；
  \item[气价差] 产品售价扣除估算天然气成本后的代理指标；
  \item[地区基差] 不同地区、运输地点和交付条件造成的价格差异；
  \item[Point-in-time] 只使用在当时已经公开的数据，避免把未来信息带入历史分析。
\end{description}

\section*{主要公开来源}

\begin{itemize}
  \item U.S. Energy Information Administration：Henry Hub 天然气；
  \item World Bank Pink Sheet：全球尿素月度价格；
  \item USDA Agricultural Marketing Service：美国肥料报价；
  \item USDA NASS：种植面积；
  \item USDA Economic Research Service：农业成本与回报；
  \item U.S. Securities and Exchange Commission：CF 财务和经营披露；
  \item CF Industries 投资者关系材料：产品关系、天然气消耗和经营说明；
  \item \url{https://www.sec.gov/Archives/edgar/data/1324404/000110465924055837/tm2413105d1_ex99-1.htm}
  \item \url{https://www.sec.gov/Archives/edgar/data/1324404/000132440425000006/cf-20241231.htm}
\end{itemize}

\end{document}
